\documentclass[10pt]{letter}
\usepackage[utf8]{inputenc}
\usepackage[margin=0.85in,top=0.7in,bottom=0.7in]{geometry}
\usepackage{hyperref}
\usepackage{siunitx}
\usepackage{eurosym}
\usepackage{setspace}
\setstretch{1.05}

\longindentation=0pt

\signature{Dr.\ Jan Kren, Paul Scherrer Institute, Switzerland (\texttt{jan.kren@psi.ch}) \\[0.5em]
Dr.\ Bo\v{s}tjan Zajec, CEA Paris-Saclay, France (\texttt{bostjan.zajec@cea.fr}) \\[0.5em]
Prof. Dr.\ Iztok Tiselj, University of Ljubljana, Slovenia (\texttt{iztok.tiselj@fmf.uni-lj.si})}

\address{Jan Kren \\ Laboratory for Simulation and Modelling (LSM) \\ Paul Scherrer Institute (PSI) \\ 5232 Villigen PSI, Switzerland}

\begin{document}

\begin{letter}{Editorial Office \\ Nature Communications}

\opening{Dear Editors,}

We are pleased to submit our manuscript entitled \textbf{``The wind is always blowing somewhere in Europe: A decade of evidence for continental-scale renewable baseload''} for consideration as an Article in \emph{Nature Communications}.

Wind power is intermittent at any single location, but weather is local: a blocking high over Germany can leave turbines idle while Atlantic storms drive strong generation across Iberia. Geographic diversification should, in principle, turn individually unreliable wind into collectively dependable supply. Whether Europe's ongoing expansion is actually exploiting this potential is an open question with direct policy consequences---the EU is targeting over 1,000~GW of wind by 2050 and planning hundreds of billions in grid investment.

We address this using the longest record of actual European wind production analyzed to date: hourly generation from 29 countries over 2015--2024. Unlike the reanalysis-based studies that dominate the literature, operational data capture curtailment, maintenance outages, and evolving fleet composition. Our central finding is that geographic diversification reliably halves wind variability---a stable property of European climatology---but that Europe's expansion is largely squandering this potential by concentrating new capacity in already-correlated regions. We identify which countries offer the greatest marginal diversification value, show that adding solar substantially raises the guaranteed baseload floor while cutting storage requirements, and use copula-based methods to characterize joint extreme behavior across the continent.

The work sits at the intersection of climate science, energy systems analysis, and portfolio theory, and speaks directly to ongoing debates about interconnection investment and renewable expansion strategy. We believe its interdisciplinary scope and policy relevance make it well suited for \emph{Nature Communications}.

We confirm that this manuscript has not been published elsewhere and is not under consideration by another journal. All authors have approved the manuscript.

\closing{Sincerely,}

\end{letter}
\end{document}
